\documentclass[10pt,a4paper]{article}
\usepackage[T1]{fontenc}
\usepackage{lmodern,amsmath,amssymb,microtype}
\usepackage[margin=20mm]{geometry}
\usepackage[hidelinks]{hyperref}
\pdfinfoomitdate=1
\pdftrailerid{}
\pdfsuppressptexinfo=15
\setlength{\parindent}{0pt}
\setlength{\parskip}{5pt}
\begin{document}
\begin{center}
{\Large Phase and night diagnostics of H-alpha}\par\vspace{4pt}
\texttt{exojax-escape} --- 25 November 2026
\end{center}

\textbf{Purpose.} Compare the same fixed atmosphere with different orbital
phases and nights before attributing a remaining difference to new physics.
The archived spectra have already been divided by nightly stellar masters.
We process the model once and sample it at the same retained native pixels;
we do not divide the archived data by another master.

\textbf{First construct the full-night model.}
Let $p(\phi,v)$ be the saved, already projected planet contrast,
$s(v_*)$ the stellar reference in the same continuum units, and $e$ an exposure.
The diagnostic factor $A_b$ and velocity offset $\Delta v_b$ act before reduction:
\begin{align}
q_e(v)&=A_{b(e)}p(\phi_e,v-\Delta v_{b(e)}),\nonumber\\
z_e(v_*)&=\frac{s(v_*)+\mathcal D_{u_e}q_e(v_*)}
 {\langle s+\mathcal D_{u_e}q_e\rangle_{C_e}},\nonumber\\
M_n(v_*)&=\frac{\sum_{e\in N_n}a_e(v_*)z_e(v_*)}
 {\sum_{e\in N_n}a_e(v_*)},\qquad
r_e(v_*)=\frac{z_e(v_*)}{M_{n(e)}(v_*)}-1.\label{eq:master}
\end{align}
Here $\mathcal D$ applies the supplied Doppler shift, $C_e$ specifies the
continuum normalization, and $a_e$ supplies master weights. Crucially, $N_n$
contains \emph{all} exposures in night $n$, including phases excluded from a fit.
The example retains all 143 master exposures and selects 115 for the final bins.

\textbf{Only then select phase and night.}
The same pixel selection and inverse-variance weights are used for data and model:
\begin{equation}
\widehat f_{bnk}=
\frac{\sum_{e\in B_b\cap N_n}\sum_p w_{ep}K_{epk}
       [\mathcal I_{ep}r_e]}
     {\sum_{e\in B_b\cap N_n}\sum_p w_{ep}K_{epk}}.
\label{eq:coadd}
\end{equation}
$\mathcal I$ samples the model at the archived stellar-frame pixel;
$K$ includes the stellar mask and bin assignment under the fixed reduction
velocity. No weight means an invalid bin. The two observing nights partly
overlap in phase; the late blue side is mostly masked.

\textbf{Why an off-line correction is not automatically safe.}
Even where an exposure's own planet line is absent, its denominator $M_n$
contains the planet in other exposures. Consequently a fit to
$(\mathrm{data}-r_e)$ outside the instantaneous line can still absorb an error
in the planetary signal through this shared denominator. Constant and linear
sideband extrapolations must be tested by injection and recovery before use.
No such background subtraction is adopted in the accepted diagnostic.
For a different exposure $f$ in the same night, at fixed $z_e$ and weights,
\begin{equation}
\frac{\partial r_e(v_*)}{\partial z_f(v_*)}
=-\frac{z_e(v_*)a_f(v_*)}
 {M_n(v_*)^2\sum_{j\in N_n}a_j(v_*)},\qquad f\ne e.
\label{eq:master-response}
\end{equation}
This nonzero response is the information shared through the master.

\textbf{Code and interpretation.}
In \texttt{observation.py}, \texttt{process\_phase\_spectra} constructs
equation~\eqref{eq:master}. In \texttt{phase\_observation.py},
\texttt{sample\_stellar\_residuals} and \texttt{coadd\_residual\_pixels}
implement equation~\eqref{eq:coadd};
\texttt{examples/phase\_halpha.py} supplies the experiment.
\path{mceg/contracts/escape.phase_observation.yaml} stores the equations,
Bindings, source/evidence pins and scoped Assessments.
The factors $A_b$ scale line plus continuum contrast, not temperature or
level populations. A smaller diagonal-error objective is not a calibrated
likelihood or evidence of atmospheric variability. Physical attribution remains
unknown, separately from operator tests and earlier atmospheric-root checks.

\small References:
\href{https://arxiv.org/html/2606.23899v1#S2}{Zhang et al. (2026), sections II--III};
\href{https://doi.org/10.5281/zenodo.20707255}{Zenodo 20707255, residual spectra}.
\end{document}
